% --- UNIVERSAL PREAMBLE BLOCK ---
\documentclass[12pt, a4paper]{article}
\usepackage[a4paper, top=2.5cm, bottom=2.5cm, left=2cm, right=2cm]{geometry}

% Univerzálne nastavenie pre rôzne kompilátory (rieši chybu fontspec)
\usepackage{iftex}
\ifPDFTeX
    % Nastavenia pre pdfLaTeX
    \usepackage[T1]{fontenc}
    \usepackage[utf8]{inputenc}
    \usepackage{lmodern}
    \usepackage[slovak]{babel}
\else
    % Nastavenia pre LuaLaTeX / XeLaTeX
    \usepackage{fontspec}
    \usepackage[slovak, bidi=basic, provide=*]{babel}
    \babelprovide[import, onchar=ids fonts]{slovak}
    \babelprovide[import, onchar=ids fonts]{english}
    
    % Set default/Latin font to Sans Serif (Noto) for stability
    \babelfont{rm}{Noto Sans}
    \babelfont[slovak]{rm}{Noto Sans}
\fi

% Fix lists for non-English
\usepackage{enumitem}
\setlist[itemize]{label=-}

% --- USER PACKAGES & SETTINGS ---
\usepackage{amsmath}
\usepackage{amsfonts}
\usepackage{amssymb}
\usepackage{graphicx} 
\usepackage{float}
\usepackage{booktabs} 
\usepackage{hyperref} 
\usepackage{listings}
\usepackage{xcolor}
\usepackage{tikz}
\usetikzlibrary{shapes.geometric, arrows.meta}

% Nastavenie pre formátovanie MATLAB kódu
\definecolor{codegreen}{rgb}{0,0.6,0}
\definecolor{codegray}{rgb}{0.5,0.5,0.5}
\definecolor{codepurple}{rgb}{0.58,0,0.82}
\definecolor{backcolour}{rgb}{0.95,0.95,0.92}

\lstdefinestyle{mystyle}{
    backgroundcolor=\color{backcolour},   
    commentstyle=\color{codegreen},
    keywordstyle=\color{blue},
    numberstyle=\tiny\color{codegray},
    stringstyle=\color{codepurple},
    basicstyle=\ttfamily\footnotesize,
    breakatwhitespace=false,         
    breaklines=true,                 
    captionpos=b,                    
    keepspaces=true,                 
    numbers=left,                    
    numbersep=5pt,                  
    showspaces=false,                
    showstringspaces=false,
    showtabs=false,                  
    tabsize=2
}

\lstset{style=mystyle}

% Nastavenie hyperref
\hypersetup{
    colorlinks=true,
    linkcolor=black,
    filecolor=magenta,      
    urlcolor=blue,
    pdftitle={Dokumentácia NutriSnap},
    pdfpagemode=FullScreen,
}

\begin{document}

% --- TITULNÁ STRANA ---
\begin{titlepage}
    \centering
    \vspace*{2cm}
    
    {\Large \textbf{Technická univerzita v Košiciach}}\\
    {\Large Fakulta elektrotechniky a informatiky}
    
    \vspace{4cm}
    
    {\Huge \textbf{Cloudové technológie}}\\
    
    \vspace{1.5cm}
    
    {\huge \textbf{Zadanie č. 2}}\\
    \vspace{0.5cm}
    {\Large NutriSnap – Analýza výživových hodnôt z fotografií jedál}
    
    \vspace{4cm}
    
    {\Large \textbf{Samuel Šarkan, Roman Šafranko}}\\
    
    \vspace{4cm}
    
    {\Large Akademický rok: 2025/2026}
    
\end{titlepage}

\tableofcontents
\newpage

\section{Slovný rozbor a analýza úlohy}
Cieľom tohto projektu je vývoj webovej aplikácie \textbf{NutriSnap}, ktorá slúži na rozpoznávanie výživových hodnôt z fotografií jedál. Úloha si vyžaduje návrh a implementáciu cloudového riešenia, ktoré prijíma nahraté obrázky od používateľov, odosiela ich na spracovanie pomocou umelej inteligencie a následne vracia detailné nutričné informácie (kalórie, bielkoviny, sacharidy, tuky). 

Základom systému je rozdelenie na klientsku časť (frontend) a serverovú časť (backend). Aplikácia musí umožňovať registráciu a prihlásenie používateľov (autentifikácia), ako aj ukladanie histórie ich naskenovaných jedál do databázy. Kvôli efektivite a modernému prístupu k vývoju bolo potrebné využiť cloudové služby pre rozpoznávanie obrazu, hosting databázy a beh samotnej aplikácie.

\section{Odôvodnenie zvolených technológií}
\begin{itemize}
    \item \textbf{Frontend (React, TypeScript, Vite, Tailwind CSS)}: Tieto technológie boli zvolené pre rýchly vývoj moderného, responzívneho a interaktívneho používateľského rozhrania. TypeScript zabezpečuje typovú bezpečnosť a predchádza mnohým chybám počas vývoja. Tailwind CSS umožňuje veľmi rýchle a elegantné štýlovanie aplikácie.
    \item \textbf{Backend (Node.js a Express)}: Zvolené pre asynchrónnu architektúru a vysokú výkonnosť pri spracovaní I/O operácií. Použitie JavaScriptu/TypeScriptu na frontende aj backende zjednocuje vývojový zásobník (stack) a zrýchľuje vývoj tímu.
    \item \textbf{Databáza (PostgreSQL na AWS RDS)}: Relačná databáza bola vybraná z dôvodu potreby jasne štruktúrovaných vzťahov (napríklad vzťah 1:N medzi používateľom a jeho naskenovanými jedlami). Amazon Relational Database Service (AWS RDS) bola zvolená ako cloudová služba, pretože ponúka automatické zálohovanie, vysokú dostupnosť a bezúdržbový chod databázy.
    \item \textbf{Spracovanie obrazu (Google Cloud Vision API)}: Namiesto budovania vlastných a výpočtovo náročných AI modelov bola na identifikáciu objektov na fotografiách využitá spoľahlivá a presná služba od spoločnosti Google.
    \item \textbf{Nutričné dáta (USDA API)}: Služba amerického ministerstva poľnohospodárstva slúži na získavanie presných kalorických a výživových dát o jedlách detegovaných cez Vision API.
\end{itemize}

\section{Diagram použitých služieb a prepojenia}
Nasledujúci diagram ukazuje architektúru celého systému a tok dát medzi cloudovými službami.

\vspace{1cm}
\begin{center}
\begin{tikzpicture}[
  box/.style={draw, rounded corners, minimum width=3cm, minimum height=1.5cm, align=center, thick},
  cloud/.style={draw, ellipse, minimum width=3.5cm, minimum height=1.5cm, align=center, thick, fill=blue!5},
  arrow/.style={->, thick, >=stealth}
]

% Nodes
\node[box, fill=green!10] (client) at (0, 0) {Klientska aplikácia\\(Frontend - React)};
\node[box, fill=yellow!10] (server) at (6, 0) {Backend Server\\(Node.js / Express)};
\node[box, fill=orange!10] (db) at (6, -3) {Databáza PostgreSQL\\(AWS RDS)};
\node[cloud] (vision) at (12, 1.5) {Google Cloud\\Vision API};
\node[cloud] (usda) at (12, -1.5) {USDA Nutrition\\API};

% Arrows
\draw[arrow, <->] (client) -- node[above] {REST API} node[below] {HTTP/JSON} (server);
\draw[arrow, <->] (server) -- node[left] {SQL} (db);
\draw[arrow, <->] (server) |- node[above, near start] {Obrázok / JSON} (vision);
\draw[arrow, <->] (server) |- node[below, near start] {Dopyt / Nutričné dáta} (usda);

\end{tikzpicture}
\end{center}
\vspace{1cm}

\section{Príspevok členov tímu k funkčnosti riešenia}
Aby bola práca na projekte rozdelená rovnomerne, úlohy boli rozčlenené do dvoch paralelných logických celkov (Frontend + Deployment vs. Cloud integrácie + Databáza). Nižšie je uvedený zoznam úloh a ich realizátorov.

\subsection*{Samuel Šarkan (Frontend, Backend-kostra, Deployment)}
\begin{itemize}
    \item [$\boxtimes$] \textbf{Návrh a implementácia používateľského rozhrania:} Vytvorenie dizajnu aplikácie vo frameworku React (Vite) s využitím Tailwind CSS na štýlovanie.
    \item [$\boxtimes$] \textbf{Vývoj základného backendového API:} Návrh štruktúry a endpointov pomocou Node.js a Express.
    \item [$\boxtimes$] \textbf{Autentifikácia a autorizácia:} Vytvorenie prihlasovacieho a registračného systému využívajúceho JSON Web Tokens (JWT).
    \item [$\boxtimes$] \textbf{Konfigurácia nasadenia (Deployment):} Zabezpečenie prevádzky frontendu a backendu v produkčnom prostredí s prepojením na doménu.
\end{itemize}

\subsection*{Roman Šafranko (Databáza, Cloud APIs, AI integrácia)}
\begin{itemize}
    \item [$\boxtimes$] \textbf{Návrh relačnej schémy a AWS RDS:} Vytvorenie schémy databázy (používatelia, história jedál) a nasadenie inštancie PostgreSQL na cloude (Amazon AWS RDS).
    \item [$\boxtimes$] \textbf{Integrácia AI (Google Cloud Vision API):} Prepojenie backendu s rozhraním Google Cloud na analýzu, klasifikáciu a detekciu objektov na používateľských fotografiách.
    \item [$\boxtimes$] \textbf{Integrácia USDA API:} Skriptovanie a mapovanie dát pre konvertovanie názvov surovín (z Vision API) na ich presné kalorické hodnoty a makroživiny prostredníctvom USDA API.
    \item [$\boxtimes$] \textbf{Orchestrácia a ukladanie:} Finálne zloženie a transformácia dát z externých API na backende s ich následným bezpečným zápisom do AWS RDS databázy.
\end{itemize}

\section{Dokumentácia k používaniu aplikácie}

\subsection{Prerekvizity (Setup)}
Pre úspešné spustenie projektu v lokálnom vývojovom prostredí je potrebné mať nainštalované:
\begin{itemize}
    \item Node.js (odporúčaná verzia 18+)
    \item Správcu balíkov npm alebo bun.
    \item Funkčné kľúče k externým službám: AWS RDS prístupové údaje, Google Vision API kľúč, USDA API kľúč a náhodný reťazec pre JWT secret.
\end{itemize}

\subsection{Spustenie aplikácie}
Aplikácia sa skladá z dvoch hlavných častí, ktoré je potrebné spustiť samostatne:

\textbf{1. Konfigurácia prostredia:}
\begin{itemize}
    \item V priečinku \texttt{backend} vytvorte súbor \texttt{.env} a vložte premenné prostredia: \texttt{DATABASE\_URL}, \texttt{JWT\_SECRET}, \texttt{GOOGLE\_VISION\_API\_KEY}, \texttt{USDA\_API\_KEY} a \texttt{PORT}.
\end{itemize}

\textbf{2. Spustenie Backend servera:}
\begin{verbatim}
cd backend
npm install
npm run dev
\end{verbatim}
Server bude počúvať na nakonfigurovanom porte (zvyčajne \texttt{http://localhost:8080} alebo iný definovaný port).

\textbf{3. Spustenie Frontend klienta:}
\begin{verbatim}
cd ..
npm install
npm run dev
\end{verbatim}
Klient sa spustí a bude dostupný na adrese (napr. \texttt{http://localhost:5173}).

\subsection{Vstupy do systému}
Používateľ do systému interaguje primárne prostredníctvom používateľského rozhrania:
\begin{itemize}
    \item \textbf{Textové vstupy:} Registračné a prihlasovacie údaje (meno, email, heslo).
    \item \textbf{Obrazové vstupy (Obrázky):} Používateľ nahráva fotografiu svojho jedla vo formátoch ako JPEG, PNG, WEBP.
\end{itemize}

\subsection{Výstupy zo systému}
Po úspešnom spracovaní zadanej fotografie systém používateľovi vráti na obrazovku:
\begin{itemize}
    \item Detegovaný názov pokrmu a jeho jednotlivých zložiek.
    \item Vypočítanú celkovú kalorickú hodnotu (v kcal).
    \item Rozdelenie makroživín (množstvo bielkovín, sacharidov a tukov).
    \item Uloženie tohto jedla do zoznamu histórie pre neskoršie zobrazenie na nástenke používateľa (Dashboard).
\end{itemize}

\end{document}
